\documentclass[10pt,a4paper]{amsart}
\usepackage[margin=19mm]{geometry}
\usepackage{amsmath,amssymb,mathtools}
\usepackage[scr=rsfs,cal=euler]{mathalfa}
\usepackage{xfrac}
\usepackage[unicode,hidelinks,bookmarks=false]{hyperref}
\newcommand{\ie}{i.e.\ }
\newcommand{\cf}{cf.\ }
\newcommand{\resp}{respectively\ }
\newcommand{\Subsection}{\subsection}
\providecommand{\nobreakdash}{\nobreak\hbox{-}}
\setlength{\emergencystretch}{2em}
\multlinegap0pt
\usepackage{tikz}
\usepackage[shortlabels]{enumitem}
\usepackage{mathtools}
\usepackage{amssymb}
\usepackage{dsfont}
\newtheorem{claim}{Claim}
\newcommand{\graph}{H}
\title{\bf\Large Sabotaging Mantel's Theorem}
\author{Natalie Behague, Debsoumya Chakraborti,, Xizhi Liu}
\date{Source: arXiv:2506.23794v2 (CC BY 4.0)}
\begin{document}
\maketitle

\noindent\emph{Source and licence.} \bf\Large Sabotaging Mantel's Theorem. Version arXiv:2506.23794v2 (CC BY 4.0), distributed by arXiv under the Creative Commons Attribution 4.0 International licence, http://creativecommons.org/licenses/by/4.0/, as recorded in the arXiv licence metadata for this exact version and shown on the abstract page https://arxiv.org/abs/2506.23794v2. Full licence notice: \texttt{licenses/arxiv-2506.23794-CC-BY-4.0.txt}.\par\smallskip

\noindent\textit{Editorial scope.} Counted unit: the article's claim CLAIM:B1-size-upper-bound (source0.tex lines 305--308). The article's complete original proof of it is delivered with it (source0.tex lines 309--325, one \texttt{proof} environment of 17 lines). No mathematics is written, repaired or re-derived: the statement and the proof are the article's own lines, in the article's own notation, and no retained line is substituted or re-flowed.  No further source-local context is needed: every cross-reference of every retained line resolves inside the two delivered spans.  Nothing was omitted from any retained span, so the package carries no omission record.  No retained line contains a citation, so the wrapper carries no bibliography and no bibliography entry.  No retained line contains a figure environment, an image-inclusion command or drawing code, so no image is delivered and the package declares no asset dependency.  Editorial formatting only, in addition: an AMS \texttt{amsart} preamble that restates the article's own declarations and macros used by the retained lines, namely the environments \texttt{claim} on separate counters, together with the standard packages those lines need.  The retained environments print in this document as Claim 1 in order of appearance, whereas the article prints the same items with the numbering of its own full document.  The complete native source of the article as distributed in the arXiv e-print for this exact version is bundled unchanged as \texttt{sources/180735/source0.tex}.
    \begin{claim}\label{CLAIM:B1-size-upper-bound}
        % We have $e(\mathcal{B}_{1}) \le e(\graph) \left(\Delta(\graph) - 1\right)$.
        We have $e(\mathcal{B}_{1}) \le \mathrm{N}(S_{2}, \graph)$.
    \end{claim}

    \begin{proof}[Proof of Claim~\ref{CLAIM:B1-size-upper-bound}]
        It follows from the definition of $\mathcal{B}_{1}$ that a pair $\{u,v\} \in \binom{[n]}{2}$ is contained in $\mathcal{B}_{1}$ iff there exists a vertex $w \in [n] \setminus \{u,v\}$ such that $\{w,u\} \in \graph$ and $\{w,v\} \in \graph$. In other words, $\{u,v\}$ is contained in the neighborhood $N_{\graph}(w)$ of some vertex $w$. 
        Therefore, 
        % \begin{align*}
        %     e(\mathcal{B}_{1})
        %     \le \sum_{w\in [n]} \binom{d_{\graph}(w)}{2}
        %     & \le \frac{1}{2} \sum_{w \in [n]}d_{\graph}(w) \left(d_{\graph}(w) - 1\right) \\
        %     & \le \frac{1}{2} \sum_{w \in [n]}d_{\graph}(w) \left(\Delta(\graph) - 1\right)
        %     = e(\graph) \left(\Delta(\graph) - 1\right),
        % \end{align*}
        \begin{align*}
            e(\mathcal{B}_{1})
            \le \sum_{w\in [n]} \binom{d_{\graph}(w)}{2}
            = \mathrm{N}(S_2, \graph),
        \end{align*}
        which proves Claim~\ref{CLAIM:B1-size-upper-bound}. 
    \end{proof}%CLAIM

\end{document}
